\subsubsection{先验界}\label{subsec:elliptic-a-priori-bounds}

最大值原理可用于得到有界区域内 $Lu=f$ 的解的逐点估计. 为陈述这些界, 引入以下量:
\begin{equation}\label{eq:elliptic-pointwise-ellipticity-and-drift}
\lambda(\bm{x})=\min_{\bm{\xi}\in\mathbb{R}^n\setminus\{0\}}\dfrac{a_{ij}(\bm{x})\xi_i\xi_j}{|\bm{\xi}|^2},\quad\beta=\max_{\bm{x}\in\overline\Omega}\dfrac{|\bm{b}(\bm{x})|}{\lambda(\bm{x})}.
\end{equation}


有以下结果.

\begin{Theorem}
\label{thm:elliptic-strip-a-priori-bound}
假设 $\Omega$ 有界, 并包含在两个相距 $d$ 的平行平面之间的带状区域内. 再假设 $c\leq0$. 若 $Lu=f$, 则
\begin{equation}\label{eq:elliptic-strip-a-priori-bound}
\max_{\overline\Omega}|u|\leq\max_{\partial\Omega}|u|+C\max_{\overline\Omega}\dfrac{|f|}{\lambda}.
\end{equation}
这里 $C=\exp((\beta+1)d)-1$.
\end{Theorem}

\begin{Proof}
假设 $Lu\geq f$, 并证明
\begin{equation}\label{eq:elliptic-strip-upper-bound}
\max_{\overline\Omega}u\leq\max_{\partial\Omega}u^++C\max_{\overline\Omega}\dfrac{|f^-|}{\lambda}.
\end{equation}
分别对 $u$ 和 $-u$ 应用这一不等式, 即得定理. 由于 $\lambda$ 和 $\beta$ 在坐标系旋转下不变, 可假设 $\Omega$ 包含在带状区域 $0<x_1<d$ 内. 令 $L_0=L-c$. 对 $\alpha\geq\beta+1$, 有
\begin{equation}\label{eq:elliptic-strip-exponential-barrier}
L_0\exp(\alpha x_1)=(\alpha^2a_{11}+\alpha b_1)\exp(\alpha x_1)\geq\lambda(\alpha^2-\alpha\beta)\exp(\alpha x_1)\geq\lambda.
\end{equation}
令
\begin{equation}\label{eq:elliptic-strip-comparison-function}
v=\max_{\partial\Omega}u^++(\exp(\alpha d)-\exp(\alpha x_1))\max_{\overline\Omega}\dfrac{|f^-|}{\lambda}.
\end{equation}
显然, 在 $\partial\Omega$ 上 $v\geq u$, 且在 $\Omega$ 上 $v\geq0$. 计算得
\begin{equation}\label{eq:elliptic-strip-comparison-inequality}
L(v-u)=L_0v+cv-Lu\leq-\lambda\max_{\overline\Omega}\dfrac{|f^-|}{\lambda}-f\leq0.
\end{equation}
Corollary~\ref{cor:elliptic-comparison} 表明 $u\leq v$, 由此得到所需结果, 其中 $C=\exp(\alpha d)-1$.
\end{Proof}

以下推论表明, 在某些情形下可以去掉 $c\leq0$ 的要求.

\begin{Corollary}
\label{cor:elliptic-small-positive-zero-order}
设有界区域 $\Omega$ 内有 $Lu=f$. 取 $C$ 如 Theorem~\ref{thm:elliptic-strip-a-priori-bound} 中所述, 并假设
\begin{equation}\label{eq:elliptic-positive-zero-order-smallness}
C^*=1-C\max_{\overline\Omega}\dfrac{c^+}{\lambda}>0.
\end{equation}
那么
\begin{equation}\label{eq:elliptic-small-positive-zero-order-bound}
\max_{\overline\Omega}|u|\leq\dfrac1{C^*}\left(\max_{\partial\Omega}|u|+C\max_{\overline\Omega}\dfrac{|f|}{\lambda}\right).
\end{equation}

\end{Corollary}

对方程 $(L_0+c^-)u=f-c^+u$ 应用前一个定理的结果, 即得证明.

\paragraph*{Problems}

\begin{enumerate}[label=\thesection.\arabic*.,ref=\thesection.\arabic*,start=1,series=book-problems]

\item\label{ex:semilinear-cubic-maximum-bound} 设 $u$ 为有界区域 $\Omega$ 上 $\Delta u=u^3-u$ 的一个解. 假设在 $\partial\Omega$ 上 $u=0$. 证明整个 $\Omega$ 内都有 $u\in[-1,1]$. 能否达到值 $\pm1$?

\item\label{ex:trace-positive-semidefinite-product} 假设 $n\times n$ 矩阵 $\bm{A}$ 和 $\bm{B}$ 对称且半正定. 证明 $\operatorname{tr}(\bm{A}\bm{B})\geq0$. Hint: $\bm{B}=\sum_k\lambda_k\bm{q}_k\bm{q}_k^T$, 其中 $\lambda_k$ 和 $\bm{q}_k$ 为 $\bm{B}$ 的特征值与特征向量.

\item\label{ex:weak-maximum-positive-zero-order-counterexample} 给出反例, 说明若 $c>0$, 则 Corollary~\ref{cor:weak-maximum-nonpositive-zero-order} 不成立.

\item\label{ex:unbounded-elliptic-comparison-counterexample} 证明, 若 $\Omega$ 无界, 则 Corollary~\ref{cor:elliptic-comparison} 不成立. Hint: 当 $\Omega$ 为平行平面之间的带状区域时, 考虑问题 $\Delta u=0$, 在 $\partial\Omega$ 上 $u=0$.

\item\label{ex:interior-tangent-ball} 若 $\partial\Omega$ 属于 $C^2$ 类且 $\bm{x}_0$ 位于 $\partial\Omega$ 上, 证明存在一个完全位于 $\Omega$ 内的球, 其边界包含 $\bm{x}_0$.

\item\label{ex:drift-neumann-adjoint-solvability} (a) 在具有光滑边界的有界区域 $\Omega$ 上, 设 $u$ 为问题
\begin{equation}\label{eq:drift-neumann-problem}
\Delta u+a_i(\bm{x})\dfrac{\partial u}{\partial x_i}=f(\bm{x}),\quad\dfrac{\partial u}{\partial n}=0\text{ on }\partial\Omega.
\end{equation}
的一个解. 假设 $\Omega$ 内有 $f\geq0$. 证明 $u$ 为常数且 $f=0$.

(b) 证明, 问题 \eqref{eq:drift-neumann-problem} 有解的必要条件是
\begin{equation}\label{eq:drift-neumann-adjoint-compatibility}
\int_\Omega f(\bm{x})v(\bm{x})\,\mathrm{d}\bm{x}=0
\end{equation}
对"伴随"方程
\begin{equation}\label{eq:drift-neumann-adjoint-equation}
\Delta v-\dfrac{\partial}{\partial x_i}\left(a_i(\bm{x})v\right)=0,\quad\dfrac{\partial v}{\partial n}-a_i(\bm{x})n_iv=0\text{ on }\partial\Omega.
\end{equation}
的每个解 $v$ 都成立.

(c) 利用后续章节将发展的技术, 可证明条件 \eqref{eq:drift-neumann-adjoint-compatibility} 也是充分的, 且 \eqref{eq:drift-neumann-adjoint-equation} 的解空间是一维的. 接受这些事实, 证明 \eqref{eq:drift-neumann-adjoint-equation} 的解要么非负, 要么非正. 形如 \eqref{eq:drift-neumann-adjoint-equation} 的方程称为 Fokker-Planck 方程, 出现在统计物理中. 只有非负解才具有物理意义.

\item\label{ex:hexagon-laplace-eigenvalue-lower-bound} 设 $\Omega$ 为边长 $a$ 的正六边形. 设 $\lambda\in\mathbb{R}$, 使得方程 $\Delta u+\lambda u=0$ 配以边界条件 $u=0$ 在 $\Omega$ 内有非平凡解. 给出 $\lambda$ 的一个下界.

\end{enumerate}
